% GENERATED FILE. Edit canonical lessons, then run npm run generate:books.
% canonical-source-hash: fnv1a64:9bef25d48d59acff
% canonical-lessons: LA-C79-cibus, LA-C79-cena, LA-C79-prandium, LA-C79-ientaculum, LA-C79-caro

\chapter{Food and Meals}
\label{ch:la-food-and-meals}
\hlchaptermodality{79}

\begin{chapteropening}
\textbf{By the end of this chapter:} \emph{I can name food, dinner, lunch, breakfast and meat.}

Complete the last lesson of chapter 79: I can name food, dinner, lunch, breakfast and meat.
\end{chapteropening}

\section[cibus, cibī]{cibus, cibī — food}

\label{lesson:LA-C79-cibus}



\begin{quote}
Before the new one: say the Latin for to fear, then the Latin for to rejoice.
\end{quote}

\subsection*{You'll want to know: cibus, cibī}
\textbf{cibus, cibī} — \textquotedblleft{}food\textquotedblright{}.

\begin{cousinweb}
\textbf{cibus}, \textquotedblleft{}food\textquotedblright{}, has no sure origin; no proposal is generally accepted.
\end{cousinweb}

\begin{grammarlens}[title={What you've built}]
One more word for this chapter.
\end{grammarlens}

\subsection*{Your turn}
\begin{itemize}
\raggedright
  \item \emph{Say it:} \emph{cibus}
  \item \emph{Say it:} \emph{cibus}, once more
  \item \emph{Say it:} say \emph{cibus}
  \item \emph{Recall:} say the Latin for to fear, then the Latin for to rejoice, then say \emph{cibus} again
\end{itemize}

\begin{tcolorbox}[breakable,colback=teal!4,colframe=teal!35!black,title={Before you move on}]
What does cibus, cibī mean? (\textquotedblleft{}Food\textquotedblright{}.) Say it once more.
\end{tcolorbox}
\section[cēna, cēnae]{cēna, cēnae — dinner, the main meal}

\label{lesson:LA-C79-cena}



\begin{quote}
Before the new one: say the Latin for to rejoice, then the Latin for food.
\end{quote}

\subsection*{You'll want to know: cēna, cēnae}
\textbf{cēna, cēnae} — \textquotedblleft{}dinner, the main meal\textquotedblright{}.

\begin{cousinweb}
\textbf{cēna} is inherited. An older spelling was \emph{caesna} (Lewis \& Short), and it is usually traced through Proto-Italic \emph{*kertsnā}, \textquotedblleft{}a portion\textquotedblright{}, to Proto-Indo-European \emph{*(s)ker-}, \textquotedblleft{}to cut\textquotedblright{} (EtymDB 2.1): the meal as one's share.
\end{cousinweb}

\begin{grammarlens}[title={What you've built}]
One more word for this chapter.
\end{grammarlens}

\subsection*{Your turn}
\begin{itemize}
\raggedright
  \item \emph{Say it:} \emph{cēna}
  \item \emph{Say it:} \emph{cēna}, once more
  \item \emph{Say it:} say \emph{cēna}
  \item \emph{Recall:} say the Latin for to rejoice, then the Latin for food, then say \emph{cēna} again
\end{itemize}

\begin{tcolorbox}[breakable,colback=teal!4,colframe=teal!35!black,title={Before you move on}]
What does cēna, cēnae mean? (\textquotedblleft{}Dinner, the main meal\textquotedblright{}.) Say it once more.
\end{tcolorbox}
\section[prandium, prandiī]{prandium, prandiī — lunch}

\label{lesson:LA-C79-prandium}



\begin{quote}
Before the new one: say the Latin for food, then the Latin for dinner.
\end{quote}

\subsection*{You'll want to know: prandium, prandiī}
\textbf{prandium, prandiī} — \textquotedblleft{}lunch\textquotedblright{}.

\begin{cousinweb}
\textbf{prandium} probably began as \textquotedblleft{}the early meal\textquotedblright{}. Its first part is an old word for \textquotedblleft{}early\textquotedblright{}, which Lewis \& Short compare with Doric Greek \emph{prān}, \textquotedblleft{}early\textquotedblright{}. The second part is often taken as the root of \emph{edō}, \textquotedblleft{}to eat\textquotedblright{}, but that is not certain.
\end{cousinweb}

\begin{grammarlens}[title={What you've built}]
One more word for this chapter.
\end{grammarlens}

\subsection*{Your turn}
\begin{itemize}
\raggedright
  \item \emph{Say it:} \emph{prandium}
  \item \emph{Say it:} \emph{prandium}, once more
  \item \emph{Say it:} say \emph{prandium}
  \item \emph{Recall:} say the Latin for food, then the Latin for dinner, then say \emph{prandium} again
\end{itemize}

\begin{tcolorbox}[breakable,colback=teal!4,colframe=teal!35!black,title={Before you move on}]
What does prandium, prandiī mean? (\textquotedblleft{}Lunch\textquotedblright{}.) Say it once more.
\end{tcolorbox}
\section[ientāculum, ientāculī]{ientāculum, ientāculī — breakfast}

\label{lesson:LA-C79-ientaculum}



\begin{quote}
Before the new one: say the Latin for dinner, then the Latin for lunch.
\end{quote}

\subsection*{You'll want to know: ientāculum, ientāculī}
\textbf{ientāculum, ientāculī} — \textquotedblleft{}breakfast\textquotedblright{}.

\begin{cousinweb}
\textbf{ientāculum} is built on \emph{ientō}, \textquotedblleft{}to have breakfast\textquotedblright{} (Lewis \& Short), with \emph{-culum}, an ending for a thing. Where \emph{ientō} comes from is uncertain.
\end{cousinweb}

\begin{grammarlens}[title={What you've built}]
One more word for this chapter.
\end{grammarlens}

\subsection*{Your turn}
\begin{itemize}
\raggedright
  \item \emph{Say it:} \emph{ientāculum}
  \item \emph{Say it:} \emph{ientāculum}, once more
  \item \emph{Say it:} say \emph{ientāculum}
  \item \emph{Recall:} say the Latin for dinner, then the Latin for lunch, then say \emph{ientāculum} again
\end{itemize}

\begin{tcolorbox}[breakable,colback=teal!4,colframe=teal!35!black,title={Before you move on}]
What does ientāculum, ientāculī mean? (\textquotedblleft{}Breakfast\textquotedblright{}.) Say it once more.
\end{tcolorbox}
\section[carō, carnis]{carō, carnis — meat}

\label{lesson:LA-C79-caro}



\begin{quote}
Before the new one: say the Latin for lunch, then the Latin for breakfast.
\end{quote}

\subsection*{You'll want to know: carō, carnis}
\textbf{carō, carnis} — \textquotedblleft{}meat\textquotedblright{}. English \emph{carnivore}.

\begin{cousinweb}
\textbf{carō} is inherited, usually traced to Proto-Indo-European \emph{*(s)ker-}, \textquotedblleft{}to cut\textquotedblright{} (EtymDB 2.1): meat as a cut piece.
\end{cousinweb}

\begin{grammarlens}[title={What you've built}]
One more word for this chapter.
\end{grammarlens}

\subsection*{Your turn}
\begin{itemize}
\raggedright
  \item \emph{Say it:} \emph{carō}
  \item \emph{Say it:} \emph{carō}, once more
  \item \emph{Say it:} say \emph{carō}
  \item \emph{Recall:} say the Latin for lunch, then the Latin for breakfast, then say \emph{carō} again
\end{itemize}

\begin{tcolorbox}[breakable,colback=teal!4,colframe=teal!35!black,title={Before you move on}]
What does carō, carnis mean? (\textquotedblleft{}Meat\textquotedblright{}.) Say it once more.
\end{tcolorbox}
